% Part: normal-modal-logic
% Chapter: sequent-calculus
% Section: introduction

\documentclass[../../../include/open-logic-section]{subfiles}

\begin{document}

\olfileid[hi]{nml}{seq}{int}

\olsection{परिचय}

प्रतिज्ञप्तिक तर्कशास्त्र के सीक्वेंट कलन को ऐसे अतिरिक्त नियमों से विस्तृत
किया जा सकता है जो \iftag{prvBox}{$\Box$\iftag{prvDiamond}{ और~}{}}{}%
\iftag{prvDiamond}{$\Diamond$}{} पर कार्य करते हैं। उदाहरणार्थ,
$\Log{K}$ के लिए हमारे पास \Log{LK} के साथ ये नियम हैं:
\[
  \iftag{prvBox}{\iftag{prvDiamond}{% <> and [] primitive
    \Axiom$\Gamma \fCenter \Delta, !A$
    \RightLabel{$\Box$}
    \UnaryInf$\Box\Gamma \fCenter \Diamond\Delta, \Box!A$
    \DisplayProof  
    \qquad
    \Axiom$!A, \Gamma \fCenter \Delta$
    \RightLabel{$\Diamond$}
    \UnaryInf$\Diamond!A, \Box\Gamma \fCenter \Diamond\Delta$
    \DisplayProof
}{% only Box primitive
\Axiom$\Gamma \fCenter !A$
\RightLabel{$\Box$}
\UnaryInf$\Box\Gamma \fCenter \Box!A$
\DisplayProof
}}{% only <> primitive
  \Axiom$!A \fCenter \Delta$
  \RightLabel{$\Diamond$}
  \UnaryInf$\Diamond!A \fCenter \Diamond\Delta$
  \DisplayProof
}
\]
$\Log{K}$ के विस्तारों के लिए और अतिरिक्त नियम भी जोड़ने होते हैं।

हर मोडल तर्कशास्त्र का ऐसा सीक्वेंट कलन नहीं होता। यहाँ तक कि $\Log{S5}$,
जो अर्थवैज्ञानिक रूप से सरल है (उसे अभिगम्यता संबंधों का तनिक भी उपयोग किए
बिना परिभाषित किया जा सकता है), के लिए भी \Log{LK} से मिलने वाला ऐसा
सीक्वेंट कलन ज्ञात नहीं है जो नियम~\Cut{} के बिना पूर्ण हो। किंतु उसका एक
कर्तन-मुक्त पूर्ण \emph{हाइपरसीक्वेंट} कलन है।

\end{document}
